\section{相关技术基础}

本章介绍本课题所依赖的核心技术基础，为后续系统设计与实验分析提供必要的理论支撑。首先阐述卷积神经网络的基本原理和本课题选用的两种骨干网络架构——ResNet与EfficientNet；其次介绍用于医学图像分割的U-Net网络结构和DAGAN生成对抗网络的合成数据原理；然后说明Grad-CAM可解释性方法的技术细节；最后讨论迁移学习和数据增强策略在本课题中的具体应用方式。


\subsection{卷积神经网络}

\subsubsection{基本原理}

卷积神经网络（Convolutional Neural Network, CNN）是一类专门设计用于处理具有网格拓扑结构数据（如图像、时间序列）的深度前馈神经网络。CNN的核心思想是通过局部感受野（Local Receptive Field）、权重共享（Weight Sharing）和池化操作（Pooling）这三个关键机制，有效提取图像中的层次化特征表示。

一个典型的卷积神经网络由多种类型的层堆叠而成：（1）\textbf{卷积层}通过可学习的卷积核在输入特征图上进行滑动窗口计算，提取局部空间特征，每个卷积核对应一种特定的特征检测器（如边缘、纹理、形状等）；（2）\textbf{激活函数层}引入非线性变换，使得网络能够学习复杂的非线性映射关系，常用的激活函数包括ReLU（Rectified Linear Unit）、Leaky ReLU、GELU等；（3）\textbf{池化层}对特征图进行下采样以降低空间分辨率，增强特征的平移不变性并减少计算量；（4）\textbf{全连接层}位于网络末端，将学习到的高层语义特征映射到输出类别空间；（5）\textbf{归一化层}如批归一化（Batch Normalization），加速训练收敛过程并提高模型稳定性。

CNN的特征提取具有明显的层次性：浅层卷积核通常学习到低级视觉特征（边缘、角点、颜色斑块等），中层卷积核组合低级特征形成中级模式（纹理、局部形状等），深层卷积核则捕获高级语义概念（物体部件、完整对象等）。这种层次化的特征学习机制使得CNN在图像分类任务中表现出色。

\subsubsection{残差网络ResNet}

随着网络深度的增加，传统的卷积神经网络面临着梯度消失/爆炸和退化（Degradation）问题：更深的网络往往不仅没有带来性能提升，反而导致训练误差和测试误差同时上升。He等人\cite{he2016deep}于2015年提出的残差网络（Residual Network, ResNet）通过引入残差学习（Residual Learning）框架有效解决了这一难题。

残差块（Residual Block）的核心思想是让网络学习输入与目标输出之间的残差映射$F(x) = H(x) - x$，而非直接学习映射$H(x)$。对于一个残差块，其输出可以表示为：
\begin{equation}
    \label{eq:resnet}
    y = F(x, \{W_i\}) + x
\end{equation}
其中$x$为残差块的输入，$F(x, \{W_i\})$表示需要学习的残差映射（通常由两层或更多层卷积构成），$y$为残差块的输出。当恒等映射是最优解时（即$H(x) = x$），将$F(x)$优化至0比直接学习恒等映射要容易得多。这一简单而优雅的设计使得训练数百甚至上千层的深度网络成为可能。

ResNet18是ResNet家族中的轻量级版本，共包含18个可学习参数层（17个卷积层+1个全连接层），参数量约为1117万。其结构包括：1个初始卷积层（7$\times$7卷积核，步长为2）、1个最大池化层、8个残差块（分为4个阶段，每阶段包含2个残差块）、1个全局平均池化层和1个全连接分类层。4个阶段的通道数依次为64、128、256、512，每个阶段内部的空间分辨率减半而通道数翻倍。ResNet18凭借其适中的模型规模和良好的性能表现，被广泛用作迁移学习的基线模型，也是本课题前期实验阶段的核心骨干网络。

\subsubsection{EfficientNet}

尽管ResNet系列取得了巨大成功，但其网络架构的设计主要依赖于人工经验和试错，缺乏系统性的缩放策略。Tan和Le\cite{tan2019efficientnet}于2019年提出的EfficientNet通过复合缩放（Compound Scaling）方法，系统性地研究了网络深度（Depth）、宽度（Width）和分辨率（Resolution）三个维度之间的最优平衡关系。

传统的人工网络缩放通常只调整其中一个维度（如只增加网络层数），这种方式效率低下且容易陷入次优配置。EfficientNet提出使用一组固定的系数来联合缩放三个维度：
\begin{equation}
    \label{eq:efficientnet}
    \begin{cases}
        depth: & d = \alpha^\phi \\
        width: & w = \beta^\phi \\
        resolution: & r = \gamma^\phi
    \end{cases}
    \quad s.t. \ \alpha \cdot \beta^2 \cdot \gamma^2 \approx 2
\end{equation}
其中$\phi$为用户指定的缩放系数，$\alpha$、$\beta$、$\gamma$为通过网格搜索确定的最优常数。约束条件$\alpha \cdot \beta^2 \cdot \gamma^2 \approx 2$确保在网络规模增加约$2^\phi$倍时，总浮点运算量（FLOPs）也同步增加约$2^\phi$倍，维持了各维度间的均衡增长。

EfficientNet-B0是EfficientNet系列的基准模型，采用MobileNet风格的倒残差 inverted bottleneck 结构作为基本构建单元，结合了 squeeze-and-excitation 注意力机制来动态调整通道权重，参数量仅为约400万，显著小于ResNet18的同时在ImageNet上达到了更高的准确率（77.1\% vs 69.3\% Top-1）。这些特性使得EfficientNet-B0特别适合小样本医学图像分类任务，因此被选为本课题最终方案的骨干网络架构。


\subsection{U-Net图像分割网络}

U-Net\cite{ronneberger2015u}是由Ronneberger等人于2015年专门针对医学图像分割任务提出的全卷积神经网络架构。该网络在ISBI 2015细胞分割挑战赛中以显著优势夺冠，此后成为医学图像分割领域最广泛使用的基线架构之一。

U-Net采用对称的编码器-解码器（Encoder-Decoder）结构，整体呈"U"形，故得名：

（1）\textbf{编码器（下采样路径）}：由连续的下采样块组成，每个下采样块包含两个$3\times3$卷积层（每个卷积后接ReLU激活函数和一个$2\times2$最大池化层（步长为2）。每经过一次池化操作，特征图的空间分辨率减半，通道数翻倍，逐步提取从低级到高级的多尺度语义特征。

（2）\textbf{解码器（上采样路径）}：由连续的上采样块组成，每个上采样块包含一个$2\times2$转置卷积层（或双线性上采样+$3\times3$卷积）进行上采样恢复空间分辨率，随后接两个$3\times3$卷积层。每经过一次上采样操作，特征图的空间分辨率加倍，通道数减半。

（3）\textbf{跳跃连接（Skip Connections）}：这是U-Net最核心的设计创新。编码器中每一层的特征图通过拼接（Concatenation）操作直接传递到解码器中对应分辨率的层。跳跃连接的作用是将编码器的高分辨率位置信息（来自浅层特征）与解码器的高层语义信息（来自深层特征）进行融合，帮助解码器更精确地定位目标边界。

（4）\textbf{输出层}：在解码器的末端，使用$1\times1$卷积将特征通道数映射到类别数，配合Sigmoid激活函数（二分类）或Softmax激活函数（多分类）输出每个像素的类别概率预测。

对于医学图像分割任务，常用的损失函数为Dice Loss，定义为：
\begin{equation}
    \label{eq:dice}
    L_{Dice} = 1 - \frac{2|P \cap G| + \epsilon}{|P| + |G| + \epsilon}
\end{equation}
其中$P$表示模型预测结果，$G$表示真实标签掩码，$\epsilon$为平滑项（通常取$10^{-6}$）以防止分母为零。Dice Loss直接优化Dice相似系数（DSC），对于处理前景/背景严重不平衡的分割问题特别有效。


\subsection{DAGAN生成对抗网络}

生成对抗网络（Generative Adversarial Network, GAN）\cite{goodfellow2014generative}是一种通过对抗过程训练生成模型的框架。GAN包含两个相互博弈的神经网络：\textbf{生成器}（Generator, $G$）学习从随机噪声向量$z$（或条件向量$c$）到真实数据空间的映射，试图生成逼真的样本；\textbf{判别器}（Discriminator, $D$）则尝试区分真实样本和生成样本。两者通过最小极大博弈（Minimax Game）交替优化：
\begin{equation}
    \label{eq:gan}
    \min_G \max_D \mathbb{E}_{x \sim p_{data}(x)}[\log D(x)] + \mathbb{E}_{z \sim p_z(z)}[\log(1 - D(G(z)))]
\end{equation}
当达到纳什均衡时，生成器产生的分布趋近于真实数据分布，判别器无法区分真假样本。

面向乳腺超声图像合成的DAGAN（Deep convolutional Adversarial Generative Network）\cite{aldhabyani2020dataset}是一种条件GAN变体，其特点在于：（1）生成器和判别器均采用深度卷积架构（DCGAN风格），使用转置卷积进行逐层上采样生成图像；（2）引入条件变量$c$（即图像类别标签：良性/恶性/正常），使生成器能够根据指定类别定向生成对应类别的超声图像；（3）在判别器中采用标签嵌入（Label Embedding）策略，将类别信息融入判别过程。

DAGAN的训练流程如下：首先从真实BUSI数据集中按类别分别训练独立的GAN模型（即每个类别训练一个DAGAN）；训练完成后，利用生成器批量合成各类别的超声图像；将合成图像按一定比例混入原始训练集中，扩充训练数据的数量和多样性，特别是增加少数类（恶性、正常）的样本量，缓解类别不平衡问题。

在本课题中，DAGAN生成的合成数据被用于Stage 3及之后的实验阶段，原始训练集546张加上2730张合成图像后，总训练样本数扩展至3276张（约6倍增幅），为后续的分类性能提升提供了重要的数据支撑。


\subsection{Grad-CAM可解释性方法}

梯度加权类激活映射（Gradient-weighted Class Activation Mapping, Grad-CAM）\cite{selvaraju2017grad}是Selvaraju等人于2017年提出的一种无需修改网络结构即可为任意CNN模型生成可视化解释的方法。该方法利用目标类别相对于最后一个卷积层输出特征图的梯度信息，定位出对预测结果贡献最大的图像区域。

Grad-CAM的具体算法流程如下：

（1）\textbf{前向传播}：给定输入图像$x$，通过网络前向传播获取最后一个卷积层的特征图$A^k$（其中$k$表示第$k$个通道）以及目标类别$c$的原始得分$y^c$。

（2）\textbf{梯度计算}：计算类别得分$y^c$相对于特征图$A^k$的梯度$\frac{\partial y^c}{\partial A^k}$，该梯度编码了每个像素对类别$c$得分的重要性。

（3）\textbf{全局平均池化求权重}：对梯度特征图进行全局平均池化，得到每个通道的重要性权重$\alpha_k^c$：
\begin{equation}
    \label{eq:gradcam_alpha}
    \alpha_k^c = \frac{1}{Z}\sum_i \sum_j \frac{\partial y^c}{\partial A_{ij}^k}
\end{equation}
其中$Z$为特征图的空间尺寸（高$\times$宽），$\alpha_k^c$反映了第$k$个特征通道对目标类别$c$的相对重要性。

（4）\textbf{加权融合与ReLU激活}：将所有通道的特征图按照对应的权重$\alpha_k^c$进行线性加权求和，并通过ReLU激活函数过滤负值：
\begin{equation}
    \label{eq:gradcam_heatmap}
    L_{Grad-CAM}^c = ReLU\left(\sum_k \alpha_k^c A^k\right)
\end{equation}
ReLU操作的含义是：只保留那些对目标类别得分有正向贡献的特征（即增加该类别得分的特征），忽略可能属于其他类别或其他图像区域的特征。

（5）\textbf{上采样与可视化}：将生成的热力图$L_{Grad-CAM}^c$通过双线性插值上采样至原始输入图像的尺寸，应用颜色映射（如热力图配色方案JET）后叠加到原图上显示。

Grad-CAM的主要优点包括：适用性广，可用于任何基于CNN的架构而无需修改网络结构或重新训练；计算高效，仅需一次反向传播即可获得热力图；可视化直观，热力图的高亮区域清晰展示了模型做出判断所依据的图像区域，便于临床医生理解和验证AI的诊断逻辑。


\subsection{迁移学习策略}

迁移学习（Transfer Learning）是指将在源域（Source Domain）大规模数据集上学到的知识迁移应用到目标域（Target Domain）任务中的机器学习方法\cite{pan2009survey}。在计算机视觉领域，最常见的迁移学习范式是在ImageNet（包含1400万张标注图像、1000个类别）等大规模自然图像数据集上预训练的CNN模型为基础，通过微调（Fine-tuning）适配到特定的下游目标任务。

对于医学图像分类这类小样本场景，迁移学习之所以有效的根本原因在于：CNN底层卷积核学习到的通用视觉特征（如边缘、纹理、形状等）在不同领域的图像之间具有一定的共享性和可迁移性。即使医学图像与自然图像存在显著的视觉差异，这些底层特征的迁移仍然能够提供良好的初始化，大幅减少目标任务所需的训练数据和训练时间，降低过拟合风险。

本课题采用的迁移学习策略包括：（1）加载ImageNet预训练的骨干网络权重作为初始化；（2）替换原始的全连接分类层，将其输出维度适配为目标任务的类别数（3类）；（3）冻结部分早期卷积层的参数（可选），仅微调深层特征和分类层；或采用全量微调策略配合较小的学习率。


\subsection{数据增强策略}

\subsubsection{传统数据增强}

传统数据增强通过对原始训练图像施加一系列预定义的几何变换和颜色变换来人工扩充训练集。本课题采用的在线数据增强策略包括以下操作：

（1）\textbf{几何变换}：随机水平翻转（概率50\%），模拟病灶可能出现在乳腺左右两侧的情况；随机旋转（$\pm10^\circ$范围内），适应超声探头角度的微小差异；随机缩放（90\%--110\%范围），模拟不同距离下的成像大小变化。

（2）\textbf{颜色抖动}：随机亮度调整（$\pm20\%$）和随机对比度调整（$\pm20\%$），模拟不同设备设置和患者个体差异造成的图像亮度和对比度变化。

需要特别强调的是，数据增强操作\textbf{仅在训练阶段}应用于训练集样本。验证集和测试集仅执行标准的预处理操作（尺寸归一化、张量化、标准化），不施加任何增强变换，以确保评估结果的客观性和可比性。

\subsubsection{基于GAN的合成数据增强}

与传统增强方法只能对已有数据进行变换重组不同，基于GAN的合成增强能够从数据分布层面生成全新的样本。DAGAN针对BUSI数据集中每个类别独立训练生成模型，能够根据类别条件定向合成对应类别的乳腺超声图像。

本课题的具体做法是：（1）使用Al-Dhabyani等人提供的预训练DAGAN模型或自行在每个类别上训练DAGAN；（2）为少数类（恶性和正常）生成更多合成样本以平衡类别分布；（3）将合成图像与原始图像合并构成扩展后的训练集，在后续训练迭代中统一参与模型优化。实验结果表明，合成数据的引入显著提升了模型在少数类上的识别准确率，尤其是良性肿瘤类别的准确率从96.49\%提升至98.25\%。


\subsection{本章小结}

本章系统地介绍了本课题所涉及的核心技术基础。首先阐述了卷积神经网络的基本原理，详细分析了两种骨干网络架构——ResNet18和EfficientNet-B0的结构特点和设计动机；其次介绍了用于医学图像分割的U-Net网络的编码器-解码器架构和跳跃连接机制；然后阐述了DAGAN生成对抗网络的工作原理及其在医学图像合成中的应用方式；接着说明了Grad-CAM可解释性方法的完整算法流程；最后讨论了迁移学习和两类数据增强策略在本课题中的具体实现方式。以上技术基础为下一章的系统设计与实现奠定了理论基础。
